\begin{center}{\Large\bfseries Response to the R53 Referee}\par\vspace{6pt}
Exact Path Representations for Finite-Catalog Renewal Design\par
R54 --- September 26, 2026\end{center}

\section*{Revision Basis and Main Changes}
We thank the referee for distinguishing the selected-boundary identity from the classical optimization machinery used after it. This revision responds to the independent report on R53 at commit \texttt{fcd7b7ab1719f00e547cd29e94cbb421fab59b0a}. It starts from that review branch and retains the reviewed R53 scientific tree. It does not mistake the earlier R49 reports for the current review.

The revision adds a direct NP-completeness reduction, an exact common-cap two-command theorem, a complete history-level construction, heterogeneous-reward primal and priced identities, and a grid-free exact search with polynomial price-path node bounds. The 142-run matched-limit comparison, ten zero-width traces, 24-case screening study, and separately declared reduction challenges evaluate those additions. The full target, realization and command constraints remain intact. Screening is now a secondary preprocessing option, not the principal publication claim.

\section{Complexity Frontier: Sections 3 and 10.1}
Theorem~\ref{thm:hard54} proves NP-completeness of the rational quadratic decision problem. The reduction already uses the fixed common reward $f(x)=x$, zero intermediate-service costs, uniform probabilities, realization ceilings equal to caps, and a nonbinding budget $m=N$. Given SUBSET SUM weights, the construction installs one free lower command and one paid optional command per history. For a subset of integer weight $s$, its exact optimized value is
\[
 \zeta-\frac{|s-W|}{2K}.
\]
Thus reaching the specified threshold is equivalent to exact subset sum. The proof verifies the original realized and expected constraints, the nonbinding budget, polynomial rational encoding, and membership in NP. It does not borrow hardness from unrestricted constrained-path problems without a reduction. The separately frozen test suite verifies 60 constructed reductions, all \HardnessSubsets{} optional subsets, and all \HardnessBooks{} feasible catalog books in the smaller cases.

Theorem~\ref{thm:twocommand54} gives a complementary exact polynomial regime beyond common eligibility: under common expected caps and a common reward, some optimum uses at most two commands. Arbitrary heterogeneous ceilings and service costs remain allowed. The proof considers a selected command at the promise, a book below the promise, and selected neighbors bracketing it, including a successor above the common cap. Enumeration of singletons and pairs has quadratic catalog dependence and polynomial rational bit complexity. All 60 common-cap regression cases agree with unrestricted book enumeration; 40 have multiple eligibility classes, and independent pair-cover certificates verify every case.

Together these results establish an explicit tractability boundary. We do not claim strong NP-hardness, hardness at budget two or three, or a complete fixed-parameter classification. The original common-prefix theorem remains exact and unchanged; it is a different sufficient polynomial regime.

\section{Constructive Proof: Sections 2 and 10.2}
The article now gives the requested individual-history construction rather than leaving it implicit in a verbal exchange. Lemma~\ref{lem:nested54} proves that a history contributing to a later terminal layer has reached its lower selected endpoint once all preceding layers are full. It then supplies a greedy weighted allocation for every partial-layer mass, with explicit adjacent-lottery probabilities. Fully traversed histories arrive at their capped terminal means, and each remaining service capacity belongs to exactly one exit set.

Lemma~\ref{lem:rearrange54} states the correct converse: any feasible arc allocation can be rearranged to an implementable allocation with the same resource and \emph{no smaller} payoff. It does not assert that an unrearranged vector must itself be implementable at exactly its original payoff. Theorem~\ref{thm:path52} combines both comparisons and gives the reconstruction sequence after its statement. Probability weights, zero capacities, capped histories, terminal edges and service feasibility are explicit. The reference implementation checks 2,035 nested constructions and rearrangement comparisons against exact original allocation.

Continuity of the convex costs on their compact domains is now explicit in the general model, matching the attainment arguments already used in the proofs. Every rational quadratic instance and every stated implemented algorithm already satisfies this regularity. The amendment does not restrict or change any computational input.

\section{Broader Mathematical Class: Sections 8 and 10.6}
We chose the methodological route requested by the referee rather than invent an application calibration. The new priced identity, Theorem~\ref{thm:price54}, allows different increasing concave terminal rewards across histories. Each history contributes its own positive priced chord increments and last-eligible service support. These terms localize at selected boundaries even though different histories need not share a common marginal ordering.

Theorem~\ref{thm:packing54} further extends the primal identity. It replaces a single common edge slope by a concave within-edge terminal allocation profile. A relaxed component allocation preserves each history's total component capacity. Packing that history's own components in decreasing marginal order preserves its individual target and improves its payoff, yielding a feasible canonical policy. This proves exactness without a common cross-history chord order. Both arbitrary-resource reconstruction and equality on optimal component vectors are checked on all 2,644 enumerated fixed books, including the heterogeneous-reward cases.

The broader polynomial additive theorem uses the new kernels; the inherited uniform-grid implementation has not been falsely relabeled as a heterogeneous solver. Only the new price implementation is used for the heterogeneous algorithm comparisons. The single aggregate equality, prefix feasibility, additive charges and pre-draw service assumptions remain explicit. We do not assert a multiple-resource or nonprefix theorem. The production-module example is retained and now includes nine exact numerical configurations specified independently of renewal histories; it is a mathematical example, not an industrial deployment claim.

\section{Approximation and Exact Search: Sections 4 and 10.3}
The article retains the polynomial uniform-resource construction and substitutes the numerical grid size into its complexity bound so that cubic catalog and quadratic budget dependence are visible. It separates arithmetic complexity, bit complexity and measured runtime. The normalized-accuracy interpretation is stated explicitly, including why a near-zero or signed value does not yield a multiplicative guarantee.

The principal new mechanism is grid-free rational price branching. One forward path calculation handles all feasible anchors and has $O(kN^2+mN^2)$ direct arithmetic complexity per price. Theorem~\ref{thm:branch54} proves original-instance bounds for every complete binary cover, finite exact termination without resource limits, and a valid interval after interruption. Zero width means exact closure. Positive width is an instance-dependent error certificate, whether or not it meets the requested tolerance. A nonconcave two-book example with a positive root duality gap is included and tested; the algorithm must branch rather than assume strong duality across books.

This is an instance-dependent alternative to uniform inverse-accuracy gridding, not a claim that a new adaptive grid was benchmarked or that a $1/\varepsilon$ lower bound was proved. The SUBSET SUM construction explains why exact decisions can require resolution of an encoding-sensitive value gap while a polynomial additive scheme remains possible. The independent checker reconstructs leaf price bounds with a backward recurrence, verifies complete binary coverage, and checks the original lottery policy. It imports neither the optimizer nor an antecedent optimization module.

\section{Meaningful General-Algorithm Evaluation: Sections 5 and 10.4}
The fixed study varies histories, catalog size, budget, accuracy, charge scale and reward heterogeneity. All 142 applicable runs execute in fresh single-threaded processes with a four-second algorithm alarm and a 2-GiB address-space cap. Peak resident memory and setup time are recorded. Numerical mixed-integer calls have explicit two-second limits for each of two envelopes and remain under the overall algorithm alarm; the external process cutoff and native-solver signal delay are disclosed. This is a matched-cap comparison, not a claim of identical resident memory or timing determinism.

Table~\ref{tab:scale54} reports all new price-path cases, including larger catalogs beyond the allotted enumeration range. The six historical fine cases are rerun under the same current limits, with enumeration times alongside all other methods in Table~\ref{tab:fine54}. Table~\ref{tab:outcomes54} retains every unsuccessful run. In the publication record, the original price study has \PriceExact{} exact closures, \PriceTolerance{} tolerance completions and \PriceLimited{} limited outcomes; all \PriceChecked{} intervals are independently verified. Ten separate zero-width runs supply the complete time-to-gap traces and Table~\ref{tab:anytime54}. No old timing is substituted for a current execution.

The subsequent complexity construction led to a separate protocol addendum committed before four reduction challenge cases were run. Those cases expose node and time limits of the new method itself. They are not silently blended into the initial protocol or represented as representative operating systems. Their exact references come from integer subset-sum dynamic programming for the constructed family. Price and enumeration request exact closure, while grid and box methods request width $1/(8K)$; this unequal target is explicit, and their outcomes are not pooled as equal-accuracy completions.

The numerical solver's original primal/dual values, node counts, model sizes, residuals and envelope errors are retained. Companion Table~\ref{tab:mip54} separates actual solve time from time allowance and approximation error. Tiny signed numerical bracket reversals remain marked as floating-point effects, not converted into exact certificates. The numerical method remains competitive on several cases; no universal speed-superiority claim is made.

\section{Screening as Secondary Preprocessing: Sections 6 and 10.5}
Proposition~\ref{prop:fix54} supplies a forced-command price bound that applies to every candidate, including possible anchors. It incorporates the promise, command budget, incumbent and competing commands, unlike the retained uniform target-box envelope. Strict upper-bound exclusion removes a command from every optimum. The equality version requires a separately retained incumbent avoiding the deletion set. These are standard incumbent-based Lagrangian inferences through the new exact original-instance oracle, and the paper credits the classical mechanism.

All 24 near-threshold benchmark cases have nonuniform caps, varying promises, budgets two through four, and mixed low and high commands. Every command has recorded exact inclusion/exclusion references, exclusion decisions, possible-anchor and wholly-ineligible status, and charge slack. Median and quartile removal rates, missed strictly irrelevant commands, incumbent cost, screening cost, and exact-reference cost are reported. The forced rule removes \ForcedAnchors{} possible-anchor occurrences across the cases, but it costs more and can miss irrelevant commands. Three separate sharpness neighbors prove computationally that a command can be necessary below the envelope threshold, tied at equality and excluded above it.

The retained envelope remains sufficient, non-anchor and potentially conservative. Its sharpness statement now expressly requires a feasible singleton and pair, budget at least two, and the full-cap promise. Its independent context checker has the stated polynomial worst-case bound. The old coarse R53 transfer interval is reported as width $1/20$ and about $21.33\%$ of its exhaustive reference; it is not called an exact solution. We do not claim to have measured the mixed-integer solver's internal presolve-removal rate. The stronger forced-path test and exact relevance give the requested problem-dependent comparison without inventing inaccessible diagnostics.

\section{Novelty, Presentation, and Preservation: Sections 7, 9, 10.7 and Minor Comments}
The introduction now compares the nearest Lagrangian path-enumeration and problem-reduction methods directly, with primary references to Handler--Zang, Beasley--Christofides and Fisher. Generic dual relaxation, tree branching and incumbent-based fixing are credited rather than presented as new. The distinct claims are the original-constraint path identities, constructive recovery, heterogeneous component extension, and the proved hardness and common-cap boundary. Capacity conservation is a supporting proposition, not a competing headline contribution.

The abstract is text-only and below 200 words; it states the quadratic catalog and linear budget/history dependence of a price oracle and the effective cubic-catalog, quadratic-budget dependence of the additive fallback. The introduction is equation-free. The current readers use anonymous title pages, 11-point text, one-inch margins, one-and-one-half spacing, author--year references and numbered tables after the references. The measured build record, not a presumed page count, determines the manuscript category. Code and data access follows the last main section.

The new branch retains every inherited source, proof, numerical result, failure and referee report. Replaced root readers have exact predecessor copies. The current companion contains the full antecedent resource, common-prefix, stability, screening and production proofs; historical experiment readers and data remain separately available rather than being silently deleted or merged into current timing claims. The file-by-file preservation audit is reproducibility evidence, not an argument for novelty or acceptance. The original joint-design scope and constraints are unchanged, and the response presents new proofs and executed evidence rather than a reduced-scope proposal.
